\clearpage
\subsubsection{Evidencia complementaria y reproducción}
\paragraph{Comprobantes internos.}
Los archivos TXT de venta y devolución de la fase 3 complementan la evidencia de P17. Los anexos indican dónde consultar los originales y las fuentes.

\paragraph{Inventario, historial y operación de Outbox.}
La Figura~\ref{fig:inventario-venta} muestra los saldos inmediatamente posteriores a la venta, antes del reintegro que presenta la Figura~\ref{fig:inventario}. La Figura~\ref{fig:cambio-precio} registra la edición del catálogo, cuya preservación histórica se coteja con la Figura~\ref{fig:pedidos}. Ambas comparaciones se corroboran en P08 y P35.
\vistafig{inventario-venta-revision.jpg}{inventario-venta}{10.55cm}{Existencias P01/P02/P03 de 7/4/19 después del cobro y de reabrir SQLite.}
\clearpage
\vistafig{cambio-precio-revision.jpg}{cambio-precio}{10.55cm}{Precio vigente de P01 a USD 21.00 al reabrir la copia del escenario de fase 4.}

\clearpage
P21 comparó dos construcciones equivalentes de 1.000 filas con los cinco SKU de control. La Figura~\ref{fig:flyweight} recoge el recuento de objetos, identidad compartida y los bytes asignados en la ejecución. La memoria se mide para la construcción comparada, en el hilo que ejecuta la medición.
\vistafig{flyweight-cierre.jpg}{flyweight}{10.7cm}{Demostración Flyweight sobre 1.000 renglones de catálogo.}
La fábrica se crea para la proyección y usa el SKU como clave de los metadatos inmutables. Los datos extrínsecos conservan valores independientes por renglón. Este resultado caracteriza esa muestra; no estima consumo total de la aplicación ni del archivo SQLite.

\clearpage
P37 indujo un fallo del notificador y reapertura del mismo archivo. La Figura~\ref{fig:outbox-fallo} corresponde al primer intento. La Figura~\ref{fig:outbox-recuperado} muestra recuperación en el segundo intento. Los archivos \ruta{evidencias/03-sqlite-antes-de-reiniciar.json} y \ruta{evidencias/08-sqlite-recuperado.json} permiten cotejar la referencia de pago y la marca persistida de procesamiento.
\vistafig{outbox-fallo.jpg}{outbox-fallo}{7.9cm}{Evento pendiente con un intento y un error simulado registrado.}
\vistafig{outbox-recuperado.jpg}{outbox-recuperado}{7.9cm}{Evento recuperado al reprocesar y métrica de sesión registrada en Auditoría.}

\clearpage
La Figura~\ref{fig:metricas-panel} completa la comparación visual de P15 junto con la Figura~\ref{fig:outbox-recuperado}. Ambas se documentaron en el recorrido de fase 4. El contador de sesión expresa la actividad observada por el host en ejecución.
\vistafig{metricas-panel.jpg}{metricas-panel}{10.7cm}{Panel general mostrando el mismo conteo que Auditoría en el recorrido de reproceso.}

\clearpage
\paragraph{Validación negativa del filtro.}
En P26, el filtro incompleto \ruta{ESTADO =} devuelve un error con ubicación en la columna 9. La Figura~\ref{fig:interpreter-error} muestra el mensaje en la interfaz ejecutada. La prueba automatizada también confirma precedencia y manejo de paréntesis.
\vistafig{interpreter-error-cierre.jpg}{interpreter-error}{10.9cm}{Error de sintaxis localizado en el monitor de reportes y comprobantes.}

